\section{经济与社会：把人、地和粮变成战争能力}

战国改革的底层是农业社会的深度动员。国家需要知道土地在哪里、谁在耕作、谁该服役、粮食怎样进仓、军队怎样获得补给。铁器、牛耕、水利、市场和人口增长常被放在同一张“生产力提高”的图里，但不同地区的速度、受益者和政治后果并不相同。

春秋时，许多资源仍通过宗族、封邑和礼仪关系流动；战国国家越来越不愿只等贵族替它转交资源。县邑、户籍、赋税、军功和法律让国家可以绕过一部分旧关系，更直接地面对土地与人。对国君而言，这意味着更快的征发与奖惩；对家庭而言，这意味着一生中更多时刻被登记、核验、征调和分类。

这种变化也并不只是压迫。军功爵、开垦奖励和更明确的身份上升通道，可能让某些原本依附于旧贵族的人获得新的机会；同时，它又把更多人推入战争和劳役的连续要求。所谓“国家变强”，在账簿上表现为粮、兵、税和文书，在人的生活里则表现为迁徙、服役、土地争夺和家族命运被更直接地卷进政治。

\subsection{粮食为什么比名将更早到达战场}

长平、伊阙、灭魏、灭楚的叙事常围绕名将展开。但一场战役能拖数月甚至数年，决定它的往往是粮食是否还能送到、道路是否还能通、城中是否还能守、后方是否还能承受征发。赵在长平的困境、楚的广域治理、秦在关中积累的优势，都不能离开这些看不见的条件。

战国末年，秦统一之所以能不断把一场胜利接到下一场胜利上，也因为它能把征服地区更快地接进文书、仓储和县级行政的网络。里耶秦简保存的不是宏大战役，而是县吏的日常文书；正是这类日常，让统一不只是一面旗帜。

\begin{turningpoint}[制度得失：直接动员的力量与代价]
\term{得}：国家可以更直接地掌握人、地、粮和兵，减少旧贵族截留资源的空间，提高战争与行政的速度。

\term{失}：同一套制度也会把风险集中到中心。征发过急、法律过密、工程过重时，家庭没有多少中间层可以缓冲。秦末的崩解并非否定战国国家化，而是提醒后人：把战时动员当作常态，国家能力会反过来制造国家承受不起的压力。
\end{turningpoint}
